\section{임피던스란 무엇인가?}

\subsection{단어의 직관}

임피던스(impedance)는 무언가의 진행을 방해하거나 억제한다는 뉘앙스를 가진다.
어원적으로 영어 단어 impede에는 방해하다, 지연시키다라는 의미가 있다.
공학에서 임피던스는 단순한 방해가 아니라, 입력된 흐름이나 운동에 대해 시스템이 어떤 방식으로 저항하고 에너지를 저장하며 다시 방출하는지를 나타내는 양이다.

여기서 말하는 방해가 반드시 에너지를 없애는 작용만 뜻하지는 않는다.
저항이나 마찰은 에너지를 열로 소산시키므로 직관적으로 방해처럼 보인다.
여기서 소산(dissipation)이란 에너지가 다시 회수하기 어려운 열, 마찰 손실, 내부 손실 같은 형태로 빠져나간다는 뜻이다.
반면 인덕터, 커패시터, 질량, 스프링은 에너지를 잠시 저장했다가 다시 돌려준다.
이들은 평균적으로 에너지를 소비하지 않을 수 있지만, 입력과 출력 사이의 위상을 바꾸고 빠른 변화에 저항한다.
저장 요소가 방해처럼 느껴지는 이유는 저장 상태를 순간적으로 바꿀 수 없기 때문이다.
전력은 effort와 flow의 곱이므로, 유한한 전력으로 자기장 전류, 전기장 전압, 질량의 속도, 스프링의 변위를 무한히 빠르게 바꾸려면 불가능할 만큼 큰 전압이나 힘이 필요하다.
그래서 저장 요소는 에너지를 영원히 없애지 않더라도, 변화가 시작되는 순간에는 그 변화를 만들어 내는 데 필요한 effort를 요구한다.
따라서 임피던스는 에너지를 열로 잃게 하는 저항성 효과와, 에너지를 저장했다가 되돌려주며 변화에 맞서는 동역학적 효과를 함께 포함한다.

이 절에서는 먼저 임피던스를 여러 물리계에 공통으로 쓰기 위한 이름표를 붙인다.
왜 이런 이름표가 필요했는지는 곧 전보선과 전화선의 역사적 문제로 돌아가 확인한다.

먼저 전력 짝을 한 장으로 고정해 두자.
전압과 전류, 힘과 속도, 토크와 각속도는 분야는 다르지만 둘을 곱하면 단위 시간당 에너지 흐름이 된다.
이 짝을 알고 있으면 뒤에서 임피던스를 볼 때 무엇을 분자로 두고 무엇을 분모로 두는지 덜 흔들린다.

\begin{table}[!htbp]
  \centering
  \caption{포트에서 함께 읽는 effort와 flow}
  \small
  \setlength{\tabcolsep}{4pt}
  \begin{tabularx}{\linewidth}{@{}>{\raggedright\arraybackslash}p{0.20\linewidth}>{\raggedright\arraybackslash}p{0.23\linewidth}>{\raggedright\arraybackslash}p{0.23\linewidth}>{\raggedright\arraybackslash}X@{}}
    \toprule
    물리계 & effort & flow & 곱하면 무엇인가 \\
    \midrule
    전기 & 전압 $v$ & 전류 $i$ & 전력 $p=vi$ \\
    병진 기계 & 힘 $f$ & 속도 $\dot{x}$ & 일률 $p=f\dot{x}$ \\
    회전 기계 & 토크 $\tau$ & 각속도 $\dot{\theta}$ & 회전 일률 $p=\tau\dot{\theta}$ \\
    \bottomrule
  \end{tabularx}
\end{table}

조금 더 일반적으로 말하면, 임피던스는 어떤 포트(port)에서 effort와 flow 사이에 생기는 동역학적 관계이다 \cite{hogan1985impedancepart1}.
여기서 포트는 시스템과 바깥세계가 에너지와 힘을 주고받는 접점이라고 생각하면 된다.
에너지 기반 모델링에서 자주 쓰는 effort와 flow라는 말은 처음에는 낯설지만, 둘의 곱이 전력이 되는 짝이라는 뜻만 붙잡으면 된다.
전기 시스템에서는 전압 $v$가 effort이고 전류 $i$가 flow이다.
기계 병진 시스템에서는 힘 $f$가 effort이고 속도 $\dot{x}$가 flow이다.
회전 시스템에서는 토크 $\tau$가 effort이고 각속도 $\dot{\theta}$가 flow이다.
말로 풀면 effort는 시스템을 밀어붙이는 변수이고, flow는 실제로 흘러가거나 움직이는 변수이다.
전압은 전하를 밀어 전류를 만들고, 힘은 물체를 밀어 속도를 만든다고 말할 수 있다.
다만 이 말은 항상 effort가 원인이고 flow가 결과라는 뜻은 아니다.
전류원처럼 전류를 직접 정하는 회로도 있고, 속도 제어기처럼 속도를 직접 명령하는 기계 시스템도 있다.
그래서 더 엄밀하게는 effort와 flow를 원인과 결과라기보다, 같은 포트에서 함께 정의되는 전력 짝(power-conjugate variables)으로 보는 편이 안전하다.
이 둘을 짝으로 묶는 이유는 단순한 관례가 아니다.
전기에서는 전압과 전류를 곱하면 전력 $p=vi$가 되고, 기계에서는 힘과 속도를 곱하면 일률 $p=f\dot{x}$가 된다.
즉 effort와 flow는 에너지가 단위 시간에 얼마나 오가는지를 계산하는 자연스러운 한 쌍이다.
이렇게 effort와 flow를 나누어 보면 전기와 기계를 같은 문장으로 비교할 수 있다.
단위로 보아도 같은 생각이 드러난다.
전기 임피던스는 $\mathrm{V/A}$, 즉 전류 $1~\mathrm{A}$를 만들기 위해 필요한 전압이고, 기계 임피던스는 $\mathrm{N/(m/s)}$, 즉 속도 $1~\mathrm{m/s}$를 만들기 위해 필요한 힘이다.

숫자로 보면 더 쉽다.
순수한 점성 댐퍼처럼 힘과 속도가 바로 비례하는 단순한 상황을 떠올려 보자.
어떤 물체가 거의 일정한 속도로 움직이는 구간에서 $1~\mathrm{N}$을 가했더니 속도가 $0.1~\mathrm{m/s}$라면, 힘과 속도의 비는
\begin{equation}
  Z_m \approx \frac{1~\mathrm{N}}{0.1~\mathrm{m/s}}
  =
  10~\mathrm{N\,s/m}
\end{equation}
이다.
이는 이 시스템이 속도 $1~\mathrm{m/s}$를 만들려면 대략 $10~\mathrm{N}$ 정도의 힘을 필요로 한다는 뜻이다.
실제 질량-스프링-댐퍼 시스템에서는 주파수와 과도응답에 따라 이 비율이 달라지므로, 이 예시는 단위와 감각을 잡기 위한 첫 번째 숫자 그림으로만 보면 된다.
여기서 위치가 아니라 속도를 나누는 이유는, 전기에서 전류가 전하가 흘러가는 빠르기인 것처럼 기계에서 속도도 위치가 변하는 빠르기이기 때문이다.
즉 임피던스는 어느 위치에 있느냐보다, 얼마나 흐르고 움직이는지를 얼마나 어렵게 만드는가를 먼저 묻는다.
따라서 각 영역에서 임피던스는 다음과 같이 쓸 수 있다.
지금부터 나오는 $Z(s)$는 시간에 따라 출렁이는 파형을 매 순간 나누겠다는 뜻이 아니다.
시간함수를 직접 다 풀기 전에, 입력과 출력의 비율을 변화의 빠르기와 주파수 관점에서 붙여 놓은 이름표라고 생각하면 된다.
스칼라 선형 시불변 시스템에서 초기 저장 에너지를 0으로 두고 볼 때, 임피던스는 flow에서 effort로 가는 전달함수처럼 읽을 수 있다.
\begin{align}
  Z_e(s) &= \frac{V(s)}{I(s)}, \\
  Z_m(s) &= \frac{F(s)}{V_x(s)}=\frac{F(s)}{sX(s)}, \\
  Z_r(s) &= \frac{\mathcal{L}\{\tau(t)\}}{\Omega_\theta(s)}
  =
  \frac{\mathcal{L}\{\tau(t)\}}{s\Theta(s)}.
\end{align}
여기서 $V_x(s)$는 병진 속도 $\dot{x}(t)$의 라플라스 변환이고, $\Omega_\theta(s)=\mathcal{L}\{\dot{\theta}(t)\}$는 단일 회전축의 각속도 라플라스 변환이다.
회전 임피던스는 토크의 라플라스 변환을 각속도의 라플라스 변환으로 나눈 값이다.
여기서 기계 포트의 힘 $F(s)$는 선택한 양의 속도 방향으로 외부가 포트에 주입해야 하는 힘을 뜻한다.
수동부호약속은 이렇게 선택한 effort와 flow의 곱이 양수일 때 에너지가 포트 안쪽 요소로 들어간다고 정하는 부호 기준이다.
실제 스프링이나 댐퍼가 물체에 가하는 복원력은 $f_s=-kx$, $f_d=-b\dot{x}$처럼 반대 방향일 수 있지만, 수동부호약속에서는 $p_m=f\dot{x}>0$일 때 외부에서 기계 요소로 전력이 들어간다고 읽는다.
이 기준으로 질량-스프링-댐퍼를 쓰면 $m\ddot{x}+b\dot{x}+kx=f_{\mathrm{in}}$이고, 뒤에서 보게 될 기계 임피던스는 $Z_m(s)=F_{\mathrm{in}}(s)/V_x(s)=ms+b+k/s$가 된다.
만약 물체가 바깥에 되돌려 주는 반작용 힘을 $F_{\mathrm{react}}$로 정의하면 $F_{\mathrm{react}}=-F_{\mathrm{in}}$이므로 부호가 반대로 보일 수 있다.
아직 $s$가 낯설다면, 일단 시간에 따른 변화를 분석하기 위한 복소수 변수라고만 받아들여도 된다.
이 관계는 전달함수처럼 초기조건을 0으로 두고 입력과 출력의 비율을 볼 때의 표기이다.
나중에 안정한 시스템의 순수한 정현파 정상상태를 볼 때는 $s=\jj\omega$를 넣는다.
전기 회로에서 직류처럼 아주 느린 최종 균형을 볼 때는 과도응답이 사라진 뒤의 $s\to 0$ 근처를 보기도 한다.
즉, $s$는 느린 변화, 빠른 진동, 감쇠되는 움직임을 한 언어로 다루기 위한 표시이다.
이 정의에서 분모가 위치가 아니라 속도라는 점이 중요하다.
기계 임피던스는 힘을 위치로 나누는 값이 아니라, 힘을 흐름인 속도로 나누는 값이다.
스프링만 따로 보면 $f=kx$라서 힘과 위치가 바로 연결되지만, 전체 기계 시스템을 전기회로처럼 비교하려면 흐름 변수인 속도를 기준으로 잡아야 한다.
그래서 스프링 항은 기계 임피던스에서 $k/s$처럼 나타난다.
정현파로 손계산해도 같은 결론이 나온다.
$x(t)=X_0\cos(\omega t)$라면 스프링 힘의 진폭은 $kX_0$이고, 속도 진폭은 $\omega X_0$이다.
같은 속도 진폭을 만들기 위해 필요한 스프링 힘의 크기는 $k/\omega$에 비례하므로, 느린 운동일수록 스프링은 힘-속도 비율에서 크게 보인다.
페이저로 쓰면 스프링 임피던스가 $k/(\jj\omega)=-\jj k/\omega$가 되고, 라플라스 표기에서는 같은 구조가 $k/s$로 나타난다.
이 때문에 정적 평형을 해석할 때는 한 가지 주의가 필요하다.
정적 평형에서는 속도 $\dot{x}$가 0이므로, 힘을 속도로 나눈 기계 임피던스 자체보다 힘-위치 관계 $F/X$, 또는 그 역수인 순응성 $X/F$가 더 직접적이다.
예를 들어 스프링의 정적 강성은 $F/X=k$이고 순응성은 $X/F=1/k$이다.
반면 기계 임피던스는 힘-속도 관계이므로 스프링 항이 $k/s$로 나타나며, $s\to0$에서 매우 커지는 형태가 된다.
처음에는 이 점이 어색하지만, 전기에서 전류가 전하의 시간 변화율인 것처럼, 기계에서 속도도 위치의 시간 변화율이라고 생각하면 자연스럽다.

반대로 flow를 effort로 나눈 양은 어드미턴스(admittance)라고 부른다.
\begin{equation}
  Y(s) = \frac{1}{Z(s)}.
\end{equation}
임피던스가 흐름을 어렵게 만드는 정도라면, 어드미턴스는 흐름을 쉽게 만드는 정도를 나타낸다.
다만 모든 상황에서 역수가 잘 정의되는 것은 아니므로, 여기서는 주로 단일 입력과 단일 출력으로 표현되는 선형 시스템에서의 직관으로 이해하면 충분하다.

\subsection{임피던스를 볼 때 붙잡을 세 가지 질문}

임피던스를 처음 배울 때는 복소수나 라플라스 변환보다 먼저 다음 세 질문을 붙잡는 편이 좋다.

\begin{enumerate}
  \item 이 요소는 에너지를 없애는가, 저장하는가?
  \item 이 요소는 무엇의 변화를 싫어하는가? 전류, 전압, 속도, 위치 중 무엇인가?
  \item 주파수가 높아지거나 낮아지면 방해가 커지는가, 작아지는가?
\end{enumerate}

저항은 전류가 흐르는 동안 에너지를 열로 없앤다.
인덕터는 전류가 갑자기 바뀌는 데 맞선다.
커패시터는 전압이 갑자기 바뀌는 데 맞선다.
질량은 속도가 갑자기 바뀌는 것을 싫어하고, 스프링은 평형점에서 위치가 벗어나는 것을 싫어하며, 댐퍼는 속도가 있는 상태 자체를 줄이려 한다.
이렇게 보면 임피던스는 단순히 막는다는 말보다 더 풍부한 개념이다.
어떤 방해는 에너지를 잃게 만들고, 어떤 방해는 에너지를 잠시 저장했다가 되돌려주며, 어떤 방해는 빠른 변화에만 강하게 나타난다.

\begin{table}[!htbp]
  \centering
  \caption{임피던스 요소를 기억하는 한 장 요약}
  \small
  \setlength{\tabcolsep}{3pt}
  \begin{tabularx}{\linewidth}{@{}>{\raggedright\arraybackslash}p{0.16\linewidth}>{\raggedright\arraybackslash}p{0.17\linewidth}>{\raggedright\arraybackslash}p{0.17\linewidth}>{\raggedright\arraybackslash}X@{}}
    \toprule
    묶음 & 전기 요소 & 기계 요소 & 기억할 말 \\
    \midrule
    소산형 & 저항 $R$ & 댐퍼 $b$ & 흐름 자체를 줄인다. 전류 $i$와 속도 $\dot{x}$가 있을 때 에너지가 열이나 내부 손실로 빠져나간다. \\
    흐름 변화 저장형 & 인덕터 $L$ & 질량 $m$ & 흐름이 갑자기 바뀌는 데 맞선다. 인덕터는 전류 변화, 질량은 속도 변화를 방해한다. \\
    기준 상태 복원형 & 커패시터 $C$ (복원항 $q/C$) & 스프링 순응성 $1/k$; 복원항은 $kx$ & 저장 상태가 기준에서 벗어날수록 대응 effort가 커진다. 역할은 닮았지만 숫자 대응은 $C\leftrightarrow 1/k$, $1/C\leftrightarrow k$이다. \\
    \bottomrule
  \end{tabularx}
  \label{tab:impedance-element-memory}
\end{table}

이 표를 한 문장으로 줄이면 다음과 같다.
저항과 댐퍼는 흐름을 열로 태워 없애고, 인덕터와 질량은 흐름의 변화에 맞서며, 커패시터와 스프링은 기준에서 벗어난 저장 상태에 대응하는 effort를 만든다.
다만 커패시터와 스프링의 숫자를 비교할 때는 조심해야 한다.
역할은 서로 닮았지만 force-voltage 비유, 즉 힘을 전압에 대응시키고 속도를 전류에 대응시키는 비유에서는 부드럽게 받아주는 정도가 $C \leftrightarrow 1/k$로 대응되고, 단단하게 되돌리는 정도가 $1/C \leftrightarrow k$로 대응된다.

\subsection{왜 임피던스가 필요했는가: 전보선과 전화선의 문제}

이제 로드맵의 첫 장면, 긴 전보선과 전화선의 문제로 돌아가 보자.
이 튜토리얼은 전신선과 전화선 문제를 임피던스를 이해하기 위한 중요한 역사적 입구로 삼는다.
임피던스 개념은 이후 교류 회로, 전력 시스템, 네트워크 해석에서도 정식화되고 확장되었으므로, 이 이야기를 유일한 기원처럼 읽을 필요는 없다.
임피던스는 시간에 따라 흔들리는 전기 신호를 멀리 보내려는 고민 속에서 자란 말이다.
책상 위의 짧은 회로에서는 전압, 전류, 저항의 관계만으로도 많은 설명이 가능하다.
하지만 선이 길어지고 신호가 펄스나 목소리처럼 계속 변하면, 저항 하나만으로는 신호가 왜 늦어지고 흐려지는지 설명하기 어렵다.
그래서 19세기 후반 전신선과 전화선 이론에서, 특히 헤비사이드가 다룬 문제에서 임피던스는 추상적인 용어가 아니라, 전보와 전화가 실제로 알아들을 수 있게 도착하느냐를 가르는 실무 언어에 가까웠다.

직류 회로만 생각하면 이야기는 비교적 단순하다.
배터리, 전선, 저항이 있고 전압을 걸면 전류가 흐른다.
전류가 충분히 안정된 뒤에는 옴의 법칙 $V=IR$만으로도 많은 것을 설명할 수 있다.
저항이 크면 전류가 작아지고, 전압이 크면 전류가 커진다.
이 정도면 흐름을 방해하는 정도를 저항 하나로 표현할 수 있다.

그런데 전보와 전화는 단순한 직류가 아니다.
전보 신호는 짧은 펄스가 켜졌다 꺼지는 신호이고, 전화 신호는 사람 목소리처럼 계속 변하는 신호이다.
즉, 전류와 전압이 시간에 따라 바뀐다.
처음에는 전선을 아주 긴 저항처럼만 생각하기 쉬웠지만, 실제 긴 케이블에서는 이상한 일이 벌어졌다.
멀리 보낸 신호가 늦게 도착하고, 펄스의 모양이 퍼지고, 전화 목소리의 높은 성분과 낮은 성분이 서로 다르게 전달되어 소리가 뭉개졌다.
전선 저항만 줄이면 된다는 직류식 생각으로는 충분하지 않았던 것이다.

여기서 한 가지 중요한 관점을 붙이면 이야기가 훨씬 쉬워진다.
짧은 전보 펄스는 겉으로는 단순히 켜짐과 꺼짐처럼 보이지만, 수학적으로는 낮은 주파수 성분과 높은 주파수 성분이 섞인 신호로 볼 수 있다.
펄스의 날카로운 모서리일수록 높은 주파수 성분이 많이 필요하다.
만약 긴 선로가 높은 주파수 성분을 더 많이 약하게 만들거나, 주파수마다 도착 시간을 다르게 만들면 어떻게 될까?
원래는 각진 모양이어야 할 펄스가 둥글게 퍼지고, 점과 선의 구분이 흐려진다.
전화에서도 비슷한 일이 생긴다.
사람 목소리에서 자음처럼 날카로운 소리에는 상대적으로 높은 주파수 성분이 많이 들어 있는데, 이 성분이 더 약해지거나 늦어지면 말소리가 멀리서 뭉개져 들린다.

이유는 전선이 단순한 저항이 아니기 때문이다.
긴 전선에는 길이를 따라 저항 $R$만 있는 것이 아니라 인덕턴스 $L$, 커패시턴스 $C$, 누설을 나타내는 컨덕턴스 $G$가 분포한다.
정확히는 선로 전체의 값 하나가 아니라 단위 길이당 값 $R'~[\Omega/\mathrm{m}]$, $L'~[\mathrm{H}/\mathrm{m}]$, $C'~[\mathrm{F}/\mathrm{m}]$, $G'~[\mathrm{S}/\mathrm{m}]$로 생각한다.
여기서 프라임($'$)은 전체 케이블 하나의 값이 아니라 1 m마다 반복해서 붙어 있는 작은 효과라는 뜻이다.
길이가 길어질수록 이 작은 효과들이 계속 누적되므로, 짧은 회로에서는 작아 보이던 저장과 누설도 긴 선로에서는 신호 모양을 바꾸는 원인이 된다.
부품 하나의 저항, 인덕턴스, 커패시턴스를 한 점에 모아 놓고 보는 관점을 집중 소자 모델(lumped model)이라고 하고, 긴 선로처럼 같은 효과가 위치를 따라 계속 퍼져 있다고 보는 관점을 분포 시스템(distributed system)이라고 부른다.

$R'$는 전선 금속 자체가 전류를 열로 잃게 만드는 정도이고, $L'$는 전류 변화가 주변 자기장에 에너지를 저장하는 정도이다.
$C'$는 전선과 대지, 또는 두 도체 사이에 전하가 저장되는 정도이다.
$G'$는 전선 금속을 따라 전류가 잘 흐르는 정도가 아니라, 절연체나 공기, 유전체를 가로질러 아주 조금 새어 나가는 병렬 누설 경로를 나타낸다.
컨덕턴스(conductance)는 저항의 반대 성격을 가진 양이며, 단위는 지멘스(S)이다.
$R'$와 $G'$는 에너지가 열이나 누설로 사라지는 손실 쪽 요소이고, $L'$와 $C'$는 에너지를 자기장이나 전기장에 잠시 저장했다가 되돌려주는 저장 쪽 요소이다.
또 $R'$와 $L'$는 선을 따라 직렬로 누적되는 효과이고, $C'$와 $G'$는 두 도체 사이 또는 도체와 대지 사이로 전하가 저장되거나 새는 병렬 효과이다.
그래서 전압의 공간 변화에는 직렬 손실과 자기장 저장이, 전류의 공간 변화에는 누설과 전기장 저장이 나타난다.
전류가 변하면 전선 주변의 자기장이 변하고, 그 변화는 전류 변화를 방해하는 유도 전압을 만든다.
또 전선과 대지, 또는 서로 나란히 놓인 두 전선 사이에는 전기장이 생기며 전하가 저장된다.
결국 긴 전선은 단순한 저항 하나가 아니라, 아주 작은 인덕터와 커패시터와 저항이 끝없이 이어진 분포 시스템처럼 행동한다.

여기서 갑자기 편미분 기호가 나오지만, 핵심은 어렵지 않다.
긴 전선을 아주 짧은 조각들로 나누어 본다고 생각하자.
길이가 $\Delta x$인 작은 조각에는 선을 따라 직렬로 $R'\Delta x$와 $L'\Delta x$가 있고, 두 도체 사이에는 병렬로 $C'\Delta x$와 $G'\Delta x$가 붙어 있다고 볼 수 있다.
한 조각 안에서도 전압과 전류는 시간에 따라 변하고, 동시에 선로를 따라 조금씩 달라진다.
따라서 긴 선로를 설명하려면 시간에 따라 얼마나 변하는가와, 전선의 위치를 따라 얼마나 변하는가를 함께 적어야 한다.
아래 식의 왼쪽은 선로를 따라 한 걸음 이동했을 때 전압이나 전류가 얼마나 바뀌는지를 말한다.
오른쪽은 그 변화가 금속 저항 손실, 자기장 에너지 저장, 절연 누설, 전기장 에너지 저장 때문에 생긴다는 뜻이다.
처음 읽을 때는 이 식을 풀려고 하지 말고, 긴 전선은 한 덩어리 저항이 아니라 작은 저장소와 손실 요소가 이어진 물체라는 표시로 읽으면 충분하다.

이 생각을 수식으로 쓰면 전신방정식(telegrapher's equations)의 기본 모양이 된다.
여기서 $x$는 선로를 따라 잰 위치, $t$는 시간, $v(x,t)$는 그 위치와 시간에서의 전압, $i(x,t)$는 전류이다.
$x$의 양의 방향을 전류가 흐르는 기준 방향으로 잡고, $v$를 도체와 귀로(return) 사이 전압으로 두면 다음처럼 쓸 수 있다.
\begin{align}
  \frac{\partial v}{\partial x} &= -R'i - L'\frac{\partial i}{\partial t}, \\
  \frac{\partial i}{\partial x} &= -G'v - C'\frac{\partial v}{\partial t}.
\end{align}
첫째 식은 선로를 따라 갈수록 전압이 바뀌는 이유가 직렬 저항 손실과 직렬 인덕턴스 전압에 있다는 뜻이다.
둘째 식은 선로를 따라 갈수록 전류가 바뀌는 이유가 병렬 누설 전류와 커패시터 충전 전류에 있다는 뜻이다.
음의 부호는 위에서 정한 위치 방향과 전압 기준에 따른 부호 약속이므로, 물리적으로는 선로 조각 안에서 전압과 전류가 손실과 저장 때문에 조금씩 달라진다는 점이 중요하다.
이 식은 저항 하나로는 긴 선로의 신호 왜곡을 설명하기 어렵다는 사실을 보여주는 출발점이다.

올리버 헤비사이드(Oliver Heaviside)는 바로 이 문제를 깊게 파고든 사람이다.
그는 전보 기사로 일하며 긴 케이블의 신호 전달 문제를 직접 접했고, 이후 맥스웰의 전자기 이론을 통신선 문제에 적용했다.
헤비사이드는 긴 케이블에서 전보 펄스가 단순히 저항성 매질 속으로 스며드는 것이 아니라 전자기파처럼 전달된다고 보았고, 긴 전보선과 전화선의 왜곡 문제를 수학적으로 다루었다 \cite{heaviside_electrical_papers,donaghyspargo2018heaviside}.
Engineering and Technology History Wiki의 전기는 이 흐름을 읽기 쉬운 역사적 배경으로 정리하고 있다 \cite{ethw_heaviside}.
특히 균일하고 거의 선형인 선로가 적절히 종단되어 있다고 해도, 감쇠가 주파수별로 달라지거나 군지연(group delay)이 주파수별로 달라지면 전화 음성처럼 여러 주파수가 섞인 신호는 멀리 갈수록 원래 모양을 잃는다.
종단 임피던스가 맞지 않으면 반사(reflection)도 생겨서 신호가 더 복잡하게 흔들린다.
헤비사이드는 선로 인덕턴스를 적절히 조절하면 왜곡을 줄일 수 있음을 보였다.
여기서 인덕턴스를 조절한다는 말은 전선의 저항만 줄이는 것이 아니라, 선로의 손실과 저장의 균형을 함께 보아야 한다는 뜻으로 이해하면 된다.
물론 이것은 모든 실제 선로에서 모든 주파수 왜곡을 완벽하게 없앨 수 있다는 뜻은 아니다.

처음 들으면 이상하게 느껴질 수 있다.
인덕턴스는 전류 변화를 방해하는 요소인데, 왜 방해 요소를 더 넣으면 신호가 좋아질까?
핵심은 신호를 가장 세게 보내는 것이 아니라, 여러 주파수 성분이 서로 비슷한 규칙으로 약해지고 지연되게 만드는 것이다.
파형은 여러 주파수 성분의 합이므로, 어떤 성분만 많이 약해지거나 파형 전체가 같은 시간만큼 늦어지는 것이 아니라 일부 성분만 더 늦어지면 전체 모양이 망가진다.
반대로 모든 성분이 비슷한 비율로 약해지고, 지연도 파형 전체가 함께 이동하는 방식으로 나타나면, 멀리 가도 모양이 덜 무너진다.

현대적인 선로 이론의 말로 하면, 균일하고 선형이며 $R'$, $L'$, $C'$, $G'$가 주파수에 거의 의존하지 않는 이상화된 선로에서는 왜곡 없는 전달 조건이
\begin{equation}
  \frac{R'}{L'} = \frac{G'}{C'}
\end{equation}
로 주어진다.
이 식을 지금 유도할 필요는 없다.
또 이 식이 임피던스의 정의라는 뜻도 아니다.
지금은 손실과 저장의 비율이 맞지 않으면 파형이 흐려지고, 이 비율을 조절하면 왜곡을 줄일 수 있다는 역사적 단서로 읽으면 충분하다.
$R'/L'$는 직렬 쪽에서 저항 손실이 자기장 저장 효과에 비해 얼마나 큰지를 나타내고, $G'/C'$는 병렬 쪽에서 누설 손실이 전기장 저장 효과에 비해 얼마나 큰지를 나타낸다.
두 양의 단위는 모두 $\mathrm{s}^{-1}$이므로, 이 조건은 서로 다른 차원의 값을 억지로 맞추는 비유가 아니라 같은 종류의 시간 척도를 맞추는 식이다.
직관적으로는 직렬 쪽의 손실/저장 비율과 병렬 쪽의 손실/저장 비율이 맞을 때, 주파수 성분마다 서로 다른 방식으로 약해지거나 파형 전체가 아닌 일부 성분만 더 늦어지는 일이 줄어든다고 읽으면 된다.
오래된 전화선에서는 절연 누설 $G'$가 작아 $G'/C'$가 작고, 선로 인덕턴스가 충분하지 않으면 $R'/L'$가 상대적으로 커질 수 있다.
이때 일정 간격의 코일로 유효 $L'$를 키우면 $R'/L'$가 낮아져, 특정 음성 대역에서 두 비율이 더 가까워지는 방향으로 움직인다.
즉 인덕턴스를 추가하는 일은 방해를 무작정 키우는 일이 아니라, 손실과 저장의 균형을 맞추어 파형 왜곡을 줄이려는 설계로 이해할 수 있다.

이 지점에서 임피던스라는 말이 왜 필요했는지 조금 더 선명해진다.
저항은 직류에서 흐르기 어려운 정도를 잘 말해 주지만, 긴 선로의 교류 신호에서는 그것만으로 부족했다.
멀리 가는 목소리와 펄스는 크기만 줄어드는 것이 아니라, 주파수 성분마다 감쇠와 군지연도 달라질 수 있었기 때문이다.
그래서 헤비사이드가 1880년대에 사용한 impedance라는 말은 단순한 새 용어가 아니라, 저항만으로 설명되지 않는 유도성ㆍ용량성 효과까지 함께 보며 신호 전달을 설명하려는 언어였다 \cite{ptb_impedance_history,heaviside_electrical_papers,donaghyspargo2018heaviside}.
직류에서는 흐르기 어려운 정도를 저항 하나로 설명해도 충분한 경우가 많다.
그러나 교류에서는 신호가 얼마나 작아지는지뿐 아니라, 전압과 전류의 흔들림이 시간상 얼마나 앞서거나 뒤지는지도 함께 보아야 한다.

역사적으로는 이것을 아주 실용적인 시험으로 생각할 수 있다.
전화선 한쪽에서 일정한 높이의 순음, 예를 들어 한 가지 주파수의 삐 소리를 계속 보낸다고 하자.
다른 쪽에서 그 소리가 얼마나 작아졌는지, 그리고 얼마나 늦게 흔들리는지를 재면 그 주파수에서의 크기 변화와 지연을 얻는다.
여러 주파수에 대해 이 일을 반복하면, 복잡한 목소리나 전보 펄스가 왜 원래 모양을 잃는지 주파수별로 볼 수 있다.
페이저와 임피던스는 바로 이런 얼마나 작아졌는가와 얼마나 늦었는가를 한 주파수에서 간단히 적기 위한 언어라고 보면 된다.

이 발상은 신호 전체를 한 번에 붙잡으려 하지 않는 데서 출발한다.
전보 펄스나 목소리는 복잡한 시간파형이지만, 하나의 정현파 성분만 떼어 놓으면 질문이 단순해진다.
이 성분은 선로를 지나며 얼마나 작아졌고, 얼마나 늦게 도착했는가?
이 질문은 긴 선로를 하나의 전달 경로로 볼 때의 감쇠와 지연 문제이다.
여러 주파수 성분마다 감쇠가 다르거나 군지연이 달라지면, 같은 음성 신호 안에서도 낮은 소리와 높은 소리가 서로 다른 규칙으로 도착해 말소리가 흐려진다.
복소수 표기는 이런 크기와 위상 정보를 함께 담는 방법이다.
그중 임피던스는 한 지점 또는 한 포트에서 전압 페이저와 전류 페이저의 비, 즉 같은 주파수에서의 $Z=V/I$를 나타낸다.
따라서 선로 전체를 지난 뒤 신호가 얼마나 약해지고 늦어지는지는 전파나 전달 특성의 문제이고, 임피던스는 그 선로의 한 지점에서 전압과 전류가 어떤 비율과 위상 차이를 갖는지를 읽는 방법이라고 구분해 두면 안전하다.
각 주파수에서 이 포트 비율을 알면 회로 소자나 선로 조각을 대수식처럼 연결해 전체 전압, 전류, 전달 특성을 계산할 수 있다.
즉 $Z=V/I$는 선로 전체의 모든 현상을 혼자 설명하는 답이 아니라, 복잡한 선로를 작은 포트 관계들의 조합으로 풀기 위한 기본 부품이다.

이렇게 정의한 임피던스의 크기는 같은 전류 진폭을 만들기 위해 필요한 전압 진폭의 비율이고, 위상은 전압 위상과 전류 위상의 차이를 말한다.
이 덕분에 미분방정식으로 풀어야 하는 복잡한 교류 회로도 정현파 정상상태에서는
\begin{equation}
  V = IZ
\end{equation}
라는 옴의 법칙과 비슷한 대수식으로 다룰 수 있다.
여기서 임피던스의 장점이 드러난다.
처음에는 전선에서 신호가 왜 뭉개질까라는 매우 현실적인 문제였는데, 그 답을 찾다 보니 저항, 자기장, 전기장, 위상, 복소수, 주파수 응답이 하나의 개념으로 묶이게 된 것이다.

\subsection{페이저와 정현파 정상상태를 아주 간단히 이해하기}

앞으로 나오는 $V=IZ$는 순간순간의 전압과 전류를 그대로 곱하는 식이 아니다.
여기서 대문자 $V$와 $I$는 시간에 따라 흔들리는 순간값이 아니라, 진폭과 위상을 담은 페이저이다.
피크값을 쓰든 RMS 값을 쓰든 한 가지 약속을 일관되게 정하면 임피던스 비율은 같게 읽을 수 있다.
또 이 설명은 입력 주파수가 그대로 출력에 남는 선형 시불변 회로의 정현파 정상상태를 전제로 한다.
아직 LTI를 배우기 전이라면, 여기서 선형시불변은 같은 주파수 입력을 넣으면 충분히 시간이 지난 뒤 같은 주파수 출력이 나오고, 크기와 위상만 바뀌는 경우 정도로 읽으면 된다.
교류 신호가 충분히 오래 지나 과도응답이 사라지고, 입력과 같은 주파수로 반복되는 상태를 정현파 정상상태(sinusoidal steady state)라고 한다.
예를 들어 전압을 $\cos(\omega t)$ 모양으로 계속 걸면, 안정한 회로의 전류도 결국 같은 $\omega$로 흔들린다.
달라지는 것은 진폭과 위상이다.

페이저(phasor)는 이 진폭과 위상을 한 번에 적어두는 복소수 표기이다.
시간에 따라 계속 도는 정현파를 매번 $\cos(\omega t+\phi)$로 쓰는 대신, 크기와 각도만 가진 화살표처럼 표현하는 것이다.
이렇게 하면 미분방정식으로 풀어야 할 교류 회로가 복소수 대수식으로 바뀐다.
허수 단위 $\jj$는 단순한 장식이 아니라, $90^\circ$ 위상 차이를 표시하는 약속이다.
인덕터와 커패시터의 임피던스에 $\jj$가 들어가는 이유도 이 때문이다.

숫자로 보면 페이저가 하는 일이 더 선명하다.
예를 들어
\begin{equation}
  v(t)=10\cos(100t), \qquad i(t)=2\cos(100t-30^\circ)
\end{equation}
라면 전류의 진폭은 전압의 $1/5$이고, 전류는 전압보다 $30^\circ$ 늦다.
따라서 이 주파수에서 임피던스의 크기는 $|Z|=5~\Omega$이고 위상은 $+30^\circ$라고 말할 수 있다.
즉 $V=IZ$는 전압과 전류의 크기 비율뿐 아니라, 누가 누구보다 먼저 흔들리는지도 함께 담는다.

\subsection{저항, 리액턴스, 임피던스}

전기회로에서 가장 단순한 방해 요소는 저항(resistance)이다.
직류 회로에서는 옴의 법칙
\begin{equation}
  V = I R
\end{equation}
만으로도 전압 $V$, 전류 $I$, 저항 $R$의 관계를 설명할 수 있다.
하지만 교류 회로에서는 전류와 전압이 시간에 따라 변한다.
인덕터는 자기장에 에너지를 저장하고, 커패시터는 전기장에 에너지를 저장한다.
이 두 소자는 단순히 에너지를 소모하는 것이 아니라, 신호의 주파수에 따라 전압과 전류 사이의 위상을 바꾼다.
이때 에너지를 저장하는 요소가 만드는 주파수 의존적인 방해를 리액턴스(reactance)라고 부른다.
리액턴스의 단위도 옴($\Omega$)이지만, 저항처럼 평균적으로 에너지를 열로 태워 없앤다는 뜻은 아니다.
저항이 임피던스의 실수부처럼 에너지를 소산하는 성분이라면, 리액턴스는 임피던스의 허수부처럼 에너지를 저장했다가 되돌려주며 위상을 바꾸는 성분이다.
그래서 교류 임피던스는 흔히 $Z=R+\jj X$처럼 쓰고, 여기서 $X$가 리액턴스이다.
즉 $X=\operatorname{Im}Z$이고, 인덕터의 리액턴스는 양의 방향, 커패시터의 리액턴스는 음의 방향으로 나타난다.
여기서 양수와 음수는 좋은 방해와 나쁜 방해를 가르는 표시가 아니다.
전압과 전류 사이의 위상 방향이 서로 반대라는 표시이다.

이때 저항, 인덕턴스, 커패시턴스를 하나의 대수적 관계로 묶기 위해 임피던스가 도입된다.
조금 더 엄밀히 말하면 인덕턴스 $L$과 커패시턴스 $C$ 자체가 곧 임피던스인 것은 아니다.
정현파 주파수 $\omega$가 정해질 때 $\jj\omega L$과 $1/(\jj\omega C)$가 전압과 전류의 비율, 즉 임피던스가 된다.
따라서 이 절에서 말하는 방해는 어떤 포트와 어떤 주파수에서 단위 흐름을 만들기 위해 필요한 작용 변수의 크기와 위상을 뜻한다.
모든 상황에서 무조건 흐름을 막는다는 뜻은 아니다.
정현파 정상상태에서 저항, 인덕터, 커패시터의 임피던스는 각각 다음과 같다.
\begin{align}
  Z_R &= R, \\
  Z_L &= \jj \omega L, \\
  Z_C &= \frac{1}{\jj \omega C}=-\frac{\jj}{\omega C}.
\end{align}
여기서 $\omega$는 각주파수이고 $\omega=2\pi f$이다.
$f$가 1초에 몇 번 흔들리는지를 나타내는 주파수라면, $\omega$는 라디안 단위로 본 흔들림의 빠르기이다.
$\jj$는 허수 단위이다.
인덕터는 $\omega$가 커질수록 $\omega L$이 커지므로 빠른 전류 변화에 더 큰 전압을 요구한다.
커패시터는 $\omega$가 커질수록 $1/(\omega C)$가 작아지므로 높은 주파수 전류는 상대적으로 잘 통과시킨다.
이 말은 앞에서 말한 커패시터가 전압 변화에 맞선다는 설명과 모순이 아니다.
같은 전류를 정해 놓고 보면 높은 주파수일수록 필요한 전압이 작아지고, 반대로 전압을 갑자기 바꾸려는 관점에서는 큰 전류가 필요하다.
무엇을 입력처럼 잡고 보느냐에 따라 같은 커패시터가 다르게 설명되는 것이다.
예를 들어 같은 진폭의 $1~\mathrm{A}$ 정현파 전류를 만들고 싶다고 하자.
저항에서는 주파수와 무관하게 $R~\mathrm{V}$의 전압 진폭이 필요하다.
인덕터에서는 주파수가 높아질수록 $\omega L~\mathrm{V}$처럼 더 큰 전압이 필요하다.
반대로 커패시터에서는 같은 $1~\mathrm{A}$ 전류를 만들기 위한 전압 진폭이 $1/(\omega C)~\mathrm{V}$로 작아진다.
이렇게 읽으면 세 임피던스가 모두 전압/전류 비율이라는 점이 분명해진다.
직렬 RLC 회로의 전체 임피던스는
\begin{equation}
  Z(\jj\omega) = R + \jj \left( \omega L - \frac{1}{\omega C} \right)
\end{equation}
로 쓸 수 있다.
여기서 $\omega L$과 $1/(\omega C)$가 서로 빼지는 이유는 두 저장 요소의 위상 방향이 반대이기 때문이다.
어떤 주파수에서는 인덕터가 만드는 양의 리액턴스와 커패시터가 만드는 음의 리액턴스가 서로 상쇄된다.
이때 회로는 저항만 남은 것처럼 보일 수 있고, 이 현상이 뒤에서 공진을 이해하는 출발점이 된다.
여기서 상쇄된다는 말은 인덕터와 커패시터 내부 전압이 사라진다는 뜻이 아니다.
각 소자 안에서는 여전히 큰 전압과 에너지 교환이 생길 수 있지만, 전원에서 바라본 전체 전압 요구가 서로 반대 위상이라 합쳐져 작게 보인다는 뜻이다.
따라서 교류 회로에서는
\begin{equation}
  V = I Z
\end{equation}
라는 형태로 옴의 법칙을 확장할 수 있다.
이 식은 시간영역의 순간값 관계라기보다, 페이저 또는 주파수 영역에서의 대수적 관계로 이해해야 한다.

\subsection{힘, 흐름, 일률}

임피던스를 물리적으로 이해하려면 전력 또는 일률(power)을 같이 보는 것이 좋다.
전기 시스템의 순간 전력은
\begin{equation}
  p_e(t) = v(t)i(t)
\end{equation}
이고, 기계 병진 시스템의 순간 일률은
\begin{equation}
  p_m(t) = f(t)\dot{x}(t)
\end{equation}
이다.
둘 다 effort와 flow의 곱이다.
따라서 전압-전류 관계와 힘-속도 관계는 수학적으로 같은 구조를 가진다.
그림~\ref{fig:effort-flow-power}는 이 대응을 전기, 병진 기계, 회전 관절의 세 경우로 나누어 보여준다.

\begin{figure}[!htbp]
  \centering
  \small
  \renewcommand{\arraystretch}{1.35}
  \begin{tabularx}{\linewidth}{
      >{\raggedright\arraybackslash}p{0.17\linewidth}
      >{\raggedright\arraybackslash}p{0.18\linewidth}
      >{\raggedright\arraybackslash}p{0.18\linewidth}
      >{\raggedright\arraybackslash}p{0.21\linewidth}
      >{\raggedright\arraybackslash}X}
    \toprule
    영역 & Effort & Flow & Power & Impedance \\
    \midrule
    전기 회로
      & 전압 $v$
      & 전류 $i$
      & $p_e=vi$
      & $Z_e(s)=V(s)/I(s)$ \\
    병진 기계
      & 힘 $f$
      & 속도 $\dot{x}$
      & $p_m=f\dot{x}$
      & $Z_m(s)=F(s)/V_x(s)$ \\
    단일 회전축
      & 토크 $\tau$
      & 각속도 $\dot{\theta}$
      & $p_r=\tau\dot{\theta}$
      & $Z_r(s)=\mathcal{L}\{\tau\}/\mathcal{L}\{\dot{\theta}\}$ \\
    \bottomrule
  \end{tabularx}

  \vspace{2mm}
  \begin{tikzpicture}[
      box/.style={draw=black!40, rounded corners=2pt, fill=black!3, align=center, inner xsep=8pt, inner ysep=4pt, font=\small},
      arr/.style={-{Latex[length=2mm]}, thick, draw=black!60}
    ]
    \node[box] (effort) {밀어붙이는 양\\effort};
    \node[box, right=11mm of effort] (system) {시스템\\$Z=\mathrm{effort}/\mathrm{flow}$};
    \node[box, right=11mm of system] (flow) {실제로 흐르는 양\\flow};
    \node[box, below=6mm of system] (power) {에너지 전달률\\$p=\mathrm{effort}\times\mathrm{flow}$};
    \draw[arr] (effort) -- (system);
    \draw[arr] (system) -- (flow);
    \draw[arr] (effort.south) |- (power.west);
    \draw[arr] (flow.south) |- (power.east);
  \end{tikzpicture}

  \caption{임피던스는 선택한 포트 방향에서 effort 변수와 flow 변수의 비율이고, 두 변수의 곱은 시간영역에서 순간적인 에너지 전달률인 power가 된다. 전기 회로에서는 전압과 전류, 기계 시스템에서는 힘과 속도, 단일 회전축에서는 토크와 각속도가 같은 역할을 한다. 그림의 화살표는 설명을 위한 포트 방향이며, 항상 effort가 원인이고 flow가 결과라는 뜻은 아니다. 다자유도 로봇 관절에서는 스칼라 나눗셈보다 $p_q=\bm{\tau}^{\mathsf T}\dot{\bm q}$, $\mathcal{L}\{\bm{\tau}(t)\}=\bm{Z}_q(s)\bm{\Omega}_q(s)$처럼 벡터와 행렬 관계로 보는 편이 안전하다. 여기서 $\bm{\Omega}_q(s)=\mathcal{L}\{\dot{\bm{q}}(t)\}$이다.}
  \label{fig:effort-flow-power}
\end{figure}


저항은 전기 에너지를 열로 소산한다.
\begin{equation}
  p_R(t) = v_R(t)i(t) = R i(t)^2 \ge 0.
\end{equation}
이와 대응되는 기계 요소는 점성 감쇠기이다.
감쇠력이 운동 방향의 반대편으로 작용한다고 쓰면
\begin{equation}
  f_d(t) = -b\dot{x}(t)
\end{equation}
이다.
이때 질량이 감쇠력 때문에 잃는 일률은 힘과 속도의 곱에 마이너스를 붙여
\begin{equation}
  p_{\mathrm{diss}}(t) = -f_d(t)\dot{x}(t)=b\dot{x}(t)^2 \ge 0
\end{equation}
으로 쓴다.
즉, 저항과 감쇠기는 에너지를 저장하지 않고 소산한다.
여기서 $R i^2$와 $b\dot{x}^2$는 수동 소자가 흡수하거나 소산하는 양의 전력으로 쓴 값이다.
실제로 질량에 작용하는 감쇠력의 부호는 운동방정식에서 $-b\dot{x}$처럼 따로 정한다.

반면 인덕터와 커패시터는 에너지를 저장한다.
인덕터 전압은
\begin{equation}
  v_L(t) = L\frac{\dd i(t)}{\dd t}
\end{equation}
이고, 저장 에너지는
\begin{equation}
  E_L(t) = \frac{1}{2}L i(t)^2
\end{equation}
이다.
전류를 갑자기 바꾸려면 $\dd i/\dd t$가 커지고, 그만큼 큰 전압이 필요하다.
따라서 인덕터는 전류의 변화에 저항한다.
중요한 점은 인덕터가 큰 전류 자체를 막는다기보다, 전류가 빠르게 바뀌려면 더 큰 전압을 필요로 한다는 것이다.
큰 질량이 빠른 속도 자체보다 속도 변화, 즉 가속도에 대해 큰 힘을 요구하는 것과 닮았다.

커패시터는 전하와 전압 사이에
\begin{equation}
  q(t) = C v_C(t)
\end{equation}
의 관계를 가지며, 전류는
\begin{equation}
  i_C(t) = C\frac{\dd v_C(t)}{\dd t}
\end{equation}
이다.
저장 에너지는
\begin{equation}
  E_C(t) = \frac{1}{2}C v_C(t)^2 = \frac{q(t)^2}{2C}
\end{equation}
이다.
커패시터 전압을 갑자기 바꾸려면 큰 전류가 필요하다.
따라서 커패시터는 전압의 변화에 저항한다.

이 관점에서 방해는 두 종류로 나뉜다.
첫째, 저항과 감쇠처럼 에너지를 소산하기 때문에 흐름을 줄이는 방해가 있다.
둘째, 인덕터와 커패시터처럼 에너지를 저장하기 때문에 변화에 반응하고 위상을 지연 또는 선행시키는 방해가 있다.
임피던스는 이 두 효과를 하나의 복소수 또는 전달함수로 묶는 개념이다.

\begin{table}[!htbp]
  \centering
  \caption{전기 요소와 기계 요소의 역할 대응}
  \begin{tabularx}{\linewidth}{@{}>{\raggedright\arraybackslash}p{0.20\linewidth}>{\raggedright\arraybackslash}p{0.20\linewidth}>{\raggedright\arraybackslash}p{0.22\linewidth}>{\raggedright\arraybackslash}X@{}}
    \toprule
    역할 & 전기 요소 & 기계 요소 & 에너지/파라미터 방향 \\
    \midrule
    소산 & 저항 $R$ & 감쇠기 $b$ & $R i^2$, $b\dot{x}^2$ \\
    운동 에너지 저장 & 인덕터 $L$ & 질량 $m$ & $\frac{1}{2}Li^2$, $\frac{1}{2}m\dot{x}^2$ \\
    위치형 에너지 저장 & 커패시터 $C$ & 스프링 $k$와 순응성 $1/k$ & $C\leftrightarrow 1/k$, $1/C\leftrightarrow k$ \\
    \bottomrule
  \end{tabularx}
\end{table}

이 표는 역할의 대응을 보여준다.
다만 전기 요소의 파라미터와 기계 요소의 파라미터가 항상 같은 방향으로 대응되는 것은 아니다.
예를 들어 커패시터의 임피던스는 $1/(sC)$이고, 스프링의 기계 임피던스는 $k/s$이다.
따라서 force-voltage, 즉 impedance 비유에서는 스프링 강성 $k$가 커패시턴스 $C$ 그 자체라기보다 $1/C$와 같은 위치에 놓인다.
이 글에서는 이 force-voltage 비유만 사용한다.
force-current 또는 mobility 비유를 쓰면 기계적 effort와 flow를 전기적 전압/전류의 반대쪽 변수에 연결하므로, 전기 소자와 기계 소자의 대응표가 달라진다.
여기서는 독자의 혼란을 줄이기 위해 두 비유를 섞지 않는다.
즉, 커패시터와 스프링은 둘 다 위치형 에너지를 저장한다는 역할이 닮았지만, 저장 변수와 파라미터의 물리 단위는 다르다.
비유는 방향을 잡기 위한 도구일 뿐, 모든 단위와 숫자가 그대로 대응된다는 뜻은 아니다.

여기서 두 층을 분리하면 구분이 선명해진다.
첫 번째 층은 역할 대응이다.
커패시터와 스프링은 둘 다 에너지를 저장했다가 되돌려주는 요소라는 점에서 닮았다.
두 번째 층은 파라미터 대응이다.
커패시턴스 $C$는 같은 전압에서 더 많은 전하를 저장할 수 있는 정도이고, 스프링 순응성 $1/k$는 같은 힘에서 더 큰 변위를 허용하는 정도이다.
그래서 부드럽게 받아주는 정도만 비교하면 $C$와 $1/k$가 닮고, 단단하게 버티는 정도를 비교하면 $1/C$와 $k$가 닮는다.

\begin{table}[!htbp]
  \centering
  \caption{커패시터와 스프링 대응에서 헷갈리기 쉬운 점}
  \begin{tabular}{p{0.28\linewidth}p{0.28\linewidth}p{0.32\linewidth}}
    \toprule
    관점 & 커패시터 & 스프링 \\
    \midrule
    물리적 역할 & 전기장에 에너지 저장 & 탄성 변형에 에너지 저장 \\
    쉽게 받아주는 정도 & $C$가 클수록 같은 전압에서 더 많은 전하를 저장함 & $1/k$가 클수록 같은 힘에서 변위가 큼 \\
    임피던스 항 & $1/(sC)$ & $k/s$ \\
    \bottomrule
  \end{tabular}
\end{table}

따라서 역할만 보면 커패시터와 스프링이 닮았고, 잘 밀리는 정도를 보면 커패시턴스 $C$와 스프링 순응성 $1/k$가 닮았다.
반대로 스프링 강성 $k$를 직접 비교하면 커패시턴스의 역수 쪽과 닮는다.
한 문장으로 줄이면, 커패시터는 전압을 잘 받아주는 정도가 $C$이고, 스프링은 힘을 잘 받아주는 정도가 $1/k$이다.
그래서 역할은 커패시터와 스프링이 닮았지만, force-voltage 비유에서 숫자의 방향은 $C \leftrightarrow 1/k$, 또는 $1/C \leftrightarrow k$로 읽는 편이 안전하다.
숫자로 보면 더 확실하다.
같은 $1~\mathrm{V}$를 걸었을 때 $C=10~\mathrm{F}$인 커패시터는 $q=10~\mathrm{C}$를 저장하지만, $C=1~\mathrm{F}$인 커패시터는 $q=1~\mathrm{C}$만 저장한다.
즉 큰 $C$는 같은 전압을 더 부드럽게 받아준다.
반대로 같은 $1~\mathrm{N}$의 힘을 스프링에 가했을 때 $k=10~\mathrm{N/m}$이면 $x=0.1~\mathrm{m}$ 밀리지만, $k=100~\mathrm{N/m}$이면 $x=0.01~\mathrm{m}$만 밀린다.
즉 작은 $k$, 또는 큰 $1/k$가 같은 힘을 더 부드럽게 받아준다.

\subsection{전기적 임피던스에서 기계적 임피던스로}

여기서 기계적 임피던스 식을 먼저 한 번 보여주는 이유는 전체 그림을 미리 잡기 위해서이다.
뒤의 LTI 시스템, RLC 회로, 질량-스프링-댐퍼 장은 이 식이 왜 이런 모양이 되는지 문법과 물리적 의미를 차례로 복원하는 과정이다.
LTI는 선형 시불변(linear time-invariant)의 줄임말이다.
거칠게 말하면 같은 조건에서 같은 입력을 주면 언제나 같은 방식으로 반응하고, 입력을 더하거나 배율을 바꾸면 출력도 그에 맞게 더해지고 배율이 바뀌는 시스템이다.
첫 독해에서는 아래 식의 완전한 유도보다 전기와 기계가 같은 구조로 이어진다는 점을 잡는 편이 중요하다.
이 장에서는 먼저 형태만 소개하고, 질량-스프링-댐퍼 장에서 각 항이 실제 운동과 어떻게 연결되는지 다시 읽을 것이다.
지금은 전기와 기계가 같은 구조로 이어진다는 정도만 붙잡으면 된다.

임피던스의 핵심은 흐름에 대한 저항이다.
전기 시스템에서 흐름은 전류이고, 그 흐름과 짝을 이루는 effort는 전압이다.
기계 시스템에서는 흐름에 해당하는 양이 속도이고, 그 흐름과 짝을 이루는 effort는 힘이다.
이 관점에서 기계적 임피던스는 힘과 속도의 동역학적 비율로 정의된다.
\begin{equation}
  Z_m(s) = \frac{F(s)}{V_x(s)}
  =
  \frac{F(s)}{sX(s)}.
\end{equation}
여기서 $F(s)$는 힘의 라플라스 변환, $V_x(s)=sX(s)$는 속도의 라플라스 변환이다.
앞에서 전압을 $V$로 썼기 때문에, 혼동을 피하려고 속도는 $V_x$라고 표시했다.
질량 $m$, 감쇠계수 $b$, 강성 $k$를 갖는 질량-스프링-댐퍼 시스템은
\begin{equation}
  m \ddot{x}(t) + b \dot{x}(t) + k x(t) = f(t)
\end{equation}
로 표현된다.
영 초기조건, 즉 계산을 시작하는 순간의 위치, 속도, 저장 에너지 같은 초기 상태를 0으로 두는 조건에서 속도의 라플라스 변환을 $V_x(s)=sX(s)$라고 두면,
\begin{equation}
  \mathcal{L}\{\ddot{x}\}=s^2X(s)=sV_x(s), \qquad X(s)=\frac{V_x(s)}{s}
\end{equation}
이다.
라플라스 변환을 적용하면
\begin{equation}
  F(s) = \left(ms + b + \frac{k}{s}\right) V_x(s)
\end{equation}
이므로 기계적 임피던스는
\begin{equation}
  Z_m(s) = ms + b + \frac{k}{s}
\end{equation}
가 된다.
전기적 RLC 회로의 임피던스와 매우 유사한 구조이다.
질량은 인덕턴스처럼 운동 에너지를 저장하고, 스프링은 커패시터처럼 위치에 따른 퍼텐셜 에너지를 저장하며, 감쇠기는 저항처럼 에너지를 소산한다.
이 식을 속도 기준으로 읽으면 물리적 의미가 재미있다.
같은 속도 진폭으로 시스템을 흔든다고 생각해 보자.
빠르게 흔들면 같은 속도라도 속도 방향이 금방 바뀌므로 가속도가 커지고, 따라서 질량항이 큰 힘을 요구한다.
반대로 아주 느리게 흔들면 같은 속도를 오래 유지하는 동안 위치 변화 폭이 커진다.
위치가 많이 벗어나면 스프링이 많이 늘어나거나 줄어드므로 스프링항이 큰 힘을 요구한다.

주파수 영역에서 식을 쓰면 이 대응이 더 뚜렷하다.
\begin{equation}
  Z_m(\jj\omega) = b + \jj\left(\omega m - \frac{k}{\omega}\right).
\end{equation}
실수부 $b$는 에너지를 소산하는 감쇠항이다.
허수부에서 질량은 $+\jj\omega m$ 방향으로 나타나고, 스프링은 $-\jj k/\omega$ 방향으로 나타난다.
둘의 부호가 반대이기 때문에 어떤 주파수에서는 서로 상쇄될 수 있다.
단위도 맞다.
$\omega m$은 $(1/\mathrm{s})\mathrm{kg}=\mathrm{kg/s}$이고, $b$는 $\mathrm{N\,s/m}=\mathrm{kg/s}$이며, $k/\omega$도 $(\mathrm{N/m})/(1/\mathrm{s})=\mathrm{kg/s}$이다.
즉 세 항은 모두 힘을 속도로 나눈 기계적 임피던스 단위로 더해질 수 있다.
낮은 주파수에서는 $k/\omega$가 커져 스프링이 지배적이고, 높은 주파수에서는 $\omega m$이 커져 질량이 지배적이다.
예를 들어 $m=1~\mathrm{kg}$, $k=100~\mathrm{N/m}$라고 하자.
$\omega=1~\mathrm{rad/s}$에서는 스프링 쪽 크기 $k/\omega=100$이 질량 쪽 크기 $\omega m=1$보다 훨씬 크다.
반대로 $\omega=100~\mathrm{rad/s}$에서는 질량 쪽 크기 $\omega m=100$이 스프링 쪽 크기 $k/\omega=1$보다 훨씬 크다.
감쇠가 작고 두 항이 서로 상쇄되는 주파수 근처에서는 작은 힘으로도 큰 운동이 나타날 수 있는데, 이것이 공진(resonance)의 핵심 직관이다.

이 장을 지나면 다음 네 가지를 말할 수 있으면 충분하다.
\begin{itemize}
  \item 임피던스는 포트에서 effort와 flow 사이의 동역학적 관계이다.
  \item 저항과 댐퍼는 에너지를 소산하고, 인덕터와 질량은 흐름 변화에 맞서며, 커패시터와 스프링은 기준 상태에서 벗어난 저장량에 대응하는 effort를 만든다.
  \item 기계 임피던스는 힘을 위치가 아니라 속도로 나눈 관계이다. 정지한 최종 평형에서는 속도가 0이므로 힘-위치 관계 $F/X=k$ 또는 순응성 $X/F=1/k$를 따로 읽는다.
  \item 스프링 항이 기계 임피던스에서 $k/s$ 또는 주파수 영역에서 $-\jj k/\omega$처럼 보이고 그 크기가 $k/\omega$가 되는 이유는, 임피던스가 위치가 아니라 속도 흐름을 기준으로 힘을 나누기 때문이다.
\end{itemize}

전보선 이야기가 남긴 질문은 시간에 따라 변하는 입력이 시스템을 지나며 크기와 위상이 어떻게 바뀌는가이다.
다음 장의 LTI 시스템과 전달함수는 바로 이 질문을 입력-출력 관계로 적는 최소 문법이다.
